\clearpage\thispagestyle{otFront}\phantomsection\label{references}\pdfbookmark[0]{Sources et édition}{bm-references}\otPageTitre{Sources et édition}{Références officielles · France}\fontsize{11}{15}\selectfont
\par\medskip\textbf{1. Arrêté et calendrier}\par BO n° 10 du 5 mars 2026, arrêté du 18 février 2026 (NOR MENE2602912A), article 3 : application en Cinquième à la rentrée 2026-2027.\par\href{https://www.education.gouv.fr/bo/2026/Hebdo10/MENE2602912A}{\textcolor{otGreenDark}{Consulter la source officielle}}\par
\par\medskip\textbf{2. Programme officiel de mathématiques}\par Annexe 2 du BO : rubriques de Cinquième, automatismes et objectifs en nombres, algèbre, géométrie, données, proportionnalité et pensée informatique.\par\href{https://www.education.gouv.fr/sites/default/files/document/Annexe\%202\%20\%E2\%80\%93\%20Programme\%20de\%20math\%C3\%A9matiques\%20pour\%20le\%20cycle\%204-480716.pdf}{\textcolor{otGreenDark}{Consulter la source officielle}}\par
\par\medskip\textbf{3. Exemples pour la mise en œuvre}\par Éduscol, mars 2026 : « Mathématiques. Exemples pour la mise en œuvre des programmes 2026. Cycle 4 », exemples de réussite et conseils pour chaque objectif de Cinquième.\par\href{https://eduscol.education.gouv.fr/sites/default/files/document/exemple-mise-en-oeuvre-c4-mathematiques0pdf-126080.pdf}{\textcolor{otGreenDark}{Consulter la source officielle}}\par
\par\medskip\textbf{4. Ressources d’accompagnement}\par Éduscol : page des ressources d’accompagnement du programme de mathématiques du cycle 4.\par\href{https://eduscol.education.gouv.fr/5736/ressources-d-accompagnement-du-programme-de-mathematiques-au-cycle-4}{\textcolor{otGreenDark}{Consulter la source officielle}}\par
\par\bigskip Sources consultées le 22 octobre 2026 ; audit de conformité
au programme et aux exemples de mise en œuvre le 23 octobre 2026.
Les cours sont rédigés pour ce livre ; les exercices, défis et bilans
sont originaux.\par\medskip
Périmètre de Cinquième : sommes et différences de relatifs et de fractions ;
valeur absolue vue comme distance à zéro, sans notation exigible ;
carrés et cubes ; distributivité simple avec lettres ; équations $x+b=c$
et $ax=c$ par opérations inverses ; dépendance de grandeurs ; moyenne
simple ; probabilités en équiprobabilité ; données saisies lues et boucles
inconditionnelles simples. Les cours signalent les propriétés admises
en Cinquième : propriétés du demi-tour, réciproque des angles alternes
internes et correspondants, concours des hauteurs et des médianes, cercle
circonscrit du triangle rectangle.\par\medskip
Les nombres premiers sont présentés comme prolongement facultatif,
conformément à la rubrique « Prolongements possibles » du programme. Les défis
approfondissent les outils du cours et ne créent pas d’exigences d’examen.
L’ordre des chapitres est une proposition de progression.

\clearpage
